% FAC-EN-U0066
% Source: fac_body.tex lines 2486--2523.
% Authority: PDF physical pp. 66--67; printed pp. 261--262
% (Chapter III section 4 opening and no. 68 complete).
% French control: T0039 repairs two inherited epsilon commands to the
% printed membership sign. No printed-source correction is admitted here.

\phantomsection\label{fac:III:4}
\section*{\S4. Relations with the functors $\operatorname{Ext}_S^q$}
\addcontentsline{toc}{section}{Chapter III, \S4. Relations with the functors
$\operatorname{Ext}_S^q$}

\phantomsection\label{fac:III:4:no:68}
\textbf{68. The functors $\operatorname{Ext}_S^q$.}
We retain the notation of \hyperref[fac:III:2:no:56]{no.~56}. If $M$ and
$N$ are two graded $S$-modules, we shall denote by
$\operatorname{Hom}_S(M,N)_n$ the group of homogeneous $S$-homomorphisms
of degree $n$ from $M$ to $N$, and by $\operatorname{Hom}_S(M,N)$ the
graded group
\phantomsection\label{fac:III:4:no:68:def:graded-hom}
\[
\sum_{n\in\mathbf{Z}}\operatorname{Hom}_S(M,N)_n.
\]
It is a graded $S$-module; when $M$ is of finite type, it coincides with
the $S$-module of all $S$-homomorphisms from $M$ to $N$.

The derived functors (cf.~[6], Chapter~V) of the functor
$\operatorname{Hom}_S(M,N)$ are the functors
$\operatorname{Ext}_S^q(M,N)$, $q=0,1,\ldots$. Let us briefly recall
\hyperref[fac:III:4:no:68:construction:Ext]{their definition}.\footnote{When
$M$ is not a module of finite type, the
$\operatorname{Ext}_S^q(M,N)$ defined here may differ from the
$\operatorname{Ext}_S^q(M,N)$ defined in [6], because
$\operatorname{Hom}_S(M,N)$ does not have the same meaning in the two
cases. Nevertheless, all the proofs in [6] remain valid without change
in the case considered here; this may be seen either directly or by
applying the Appendix of [6].}

\phantomsection\label{fac:III:4:no:68:construction:Ext}
One chooses a ``resolution'' of $M$, that is, an exact sequence
\phantomsection\label{fac:III:4:no:68:formula:resolution}
\[
\cdots\longrightarrow L^{q+1}\longrightarrow L^q\longrightarrow\cdots
\longrightarrow L^0\longrightarrow M\longrightarrow0,
\]
where the $L^q$ are free graded $S$-modules and the maps are
homomorphisms (that is, as usual, homogeneous $S$-homomorphisms of
degree~$0$). Put
\phantomsection\label{fac:III:4:no:68:def:Ext}
$C^q=\operatorname{Hom}_S(L^q,N)$. The homomorphism
$L^{q+1}\longrightarrow L^q$ defines by transposition a homomorphism
$d:C^q\longrightarrow C^{q+1}$ satisfying $d\circ d=0$. Thus
\phantomsection\label{fac:III:4:no:68:formula:Hom-complex}
$C=\sum_{q\geq0}C^q$ has the structure of a complex, and its $q$-th
cohomology group is, by definition, $\operatorname{Ext}_S^q(M,N)$; one
shows that it does not depend on
\hyperref[fac:III:4:no:68:formula:resolution]{the chosen resolution}. Since
the $C^q$
are graded $S$-modules and $d:C^q\longrightarrow C^{q+1}$ is homogeneous
of degree~$0$, the $\operatorname{Ext}_S^q(M,N)$ are graded $S$-modules,
with homogeneous subspaces $\operatorname{Ext}_S^q(M,N)_n$. These latter
spaces are the cohomology groups of the complex formed by the
$\operatorname{Hom}_S(L^q,N)_n$; in other words, they are the derived
functors of $\operatorname{Hom}_S(M,N)_n$.

\phantomsection\label{fac:III:4:no:68:properties}
Let us recall the principal properties of $\operatorname{Ext}_S^q$:
$\operatorname{Ext}_S^0(M,N)=\operatorname{Hom}_S(M,N)$;
$\operatorname{Ext}_S^q(M,N)=0$ for $q>r+1$ if $M$ is of finite type
(by Hilbert's syzygy theorem, cf.~[6], Chapter~VIII, Theorem~6.5);
$\operatorname{Ext}_S^q(M,N)$ is an $S$-module of finite type if $M$ and
$N$ are of finite type (for one may choose a resolution in which the
$L^q$ are of finite type); and for every $n\in\mathbf{Z}$ there are
canonical isomorphisms
\phantomsection\label{fac:III:4:no:68:formula:shift}
\[
\operatorname{Ext}_S^q(M(n),N)
 \simeq \operatorname{Ext}_S^q(M,N(-n))
 \simeq \operatorname{Ext}_S^q(M,N)(-n).
\]

The exact sequences
\phantomsection\label{fac:III:4:no:68:formula:short-exact}
\[
0\longrightarrow N\longrightarrow N'\longrightarrow N''\longrightarrow0
\qquad\text{and}\qquad
0\longrightarrow M\longrightarrow M'\longrightarrow M''\longrightarrow0
\]
give rise to the exact sequences
\phantomsection\label{fac:III:4:no:68:formula:long-exact-N}
\[
\cdots\longrightarrow\operatorname{Ext}_S^q(M,N)
\longrightarrow\operatorname{Ext}_S^q(M,N')
\longrightarrow\operatorname{Ext}_S^q(M,N'')
\longrightarrow\operatorname{Ext}_S^{q+1}(M,N)\longrightarrow\cdots
\]
and
\phantomsection\label{fac:III:4:no:68:formula:long-exact-M}
\[
\cdots\longrightarrow\operatorname{Ext}_S^q(M'',N)
\longrightarrow\operatorname{Ext}_S^q(M',N)
\longrightarrow\operatorname{Ext}_S^q(M,N)
\longrightarrow\operatorname{Ext}_S^{q+1}(M'',N)\longrightarrow\cdots.
\]
